\chapter{基于多尺度特征融合的电机故障诊断方法}

\section{问题分析}

不同类型的电机故障在振动信号中表现出不同的尺度特征。例如，轴承局部故障产生高频冲击信号，而转子不平衡则表现为低频振动。单一尺度的特征提取难以充分描述不同故障的特征，影响诊断精度。

\section{多尺度特征融合网络设计}

\subsection{网络总体结构}

本文设计了一种多尺度特征融合网络（MSFF-Net），其核心思想是利用不同大小的卷积核提取不同尺度的特征，并通过特征融合模块将多尺度特征进行有效整合。

\subsection{多尺度特征提取模块}

多尺度特征提取模块采用并行结构，使用三种不同大小的卷积核（$K=3, 5, 7$）分别提取不同尺度的特征：

\begin{equation}
	\mathbf{h}_k^{(l)} = \text{Conv1D}_{K=k}\left(\mathbf{h}^{(l-1)}\right), \quad k \in \{3, 5, 7\}
	\label{eq:multiscale}
\end{equation}

\subsection{特征融合模块}

特征融合模块将不同尺度的特征进行拼接和加权融合：

\begin{equation}
	\mathbf{h}_{fused}^{(l)} = \sum_{k \in \{3,5,7\}} \alpha_k \cdot \mathbf{h}_k^{(l)}
	\label{eq:fusion}
\end{equation}

其中$\alpha_k$为可学习的注意力权重，满足$\sum_k \alpha_k = 1$。

\section{实验结果与分析}

\subsection{实验设置}

在CWRU数据集和江南大学轴承数据集上分别进行实验，验证MSFF-Net的有效性。

\subsection{对比实验结果}

\begin{table}[htbp]
	\centering
	\caption{不同方法在两个数据集上的诊断准确率(\%)}
	\label{tab:msff_results}
	\begin{tabular}{ccc}
		\toprule
		方法 & CWRU & 江南大学 \\
		\midrule
		标准CNN & 97.3 & 91.2 \\
		多分支CNN & 98.1 & 93.5 \\
		ResNet-1D & 97.8 & 92.8 \\
		本文MSFF-Net & 99.2 & 95.6 \\
		\bottomrule
	\end{tabular}
\end{table}

实验结果表明，本文提出的MSFF-Net在两个数据集上均取得了最优的诊断准确率，在CWRU数据集上达到99.2\%，在江南大学数据集上达到95.6\%，验证了多尺度特征融合策略的有效性。

\section{本章小结}

本章提出了基于多尺度特征融合的电机故障诊断方法，设计MSFF-Net网络结构，利用不同大小的卷积核提取多尺度特征，并通过注意力机制进行特征融合。实验结果表明，该方法在两个数据集上均取得了优于对比方法的诊断精度。